\section{Count repeated pixels independently of input bytes}
For RUN, the low six bits store the run length minus one. \texttt{c0} therefore repeats
the previous pixel once, and \texttt{fd} repeats it sixty-two times. The two remaining
full-byte values belong to RGB and RGBA. There is no sixty-three- or sixty-four-pixel
single run tag.

Compare the decoded length with the number of output pixels still required before
writing the run. A run that crosses that bound is malformed even if the destination
arena has spare mapped storage. For an accepted run, write the same color once per
pixel, advance the output pointer by four each time, and reduce the remaining pixel
count. Consume no additional input bytes for those repeats.

A long sequence needs several run chunks. Sixty-three opaque black pixels at the
start of an image can use a run of sixty-two and a run of one. The previous pixel
already supplies opaque black, so no initial color command is required. A nonblack
uniform image first establishes its color, then repeats it. At end of image, flush
any pending encoder run even when its length is less than sixty-two.

An input cursor stopped after the run tag and an output cursor partway through its
expansion are consistent states. When observing the display, count accepted output
stores rather than expecting each input byte to correspond to one visible pixel.
